\begin{frame}{Colaborando con otras personas}

    Los repositorios remotos son \textit{copias} de nuestro proyecto a las cuales accedemos a través
    de Internet. Puede haber varios, cada uno de los cuales
    puede ser de solo lectura o de lectura/escritura, según los permisos que tengamos.

    \begin{figure}[ht]
        \begin{center}
            \includegraphics[height=1.5in]{images/repo-remoto.pdf}
        \end{center}
    \end{figure}

    Colaborar con otros implica gestionar estos repositorios remotos, y mandar (\textbf{push}) y recibir (\textbf{pull})
    datos de ellos cuando necesites compartir cambios.

\end{frame}

\begin{frame}{Usando Git online}
\begin{block}{Git server as a service}
    Para subir nuestros proyectos a Internet, ya sea para colaborar con otras personas o para compartir nuestro código, podemos usar un servicio de servidor web de Git.

    Los más conocidos son Github, Codeberg y Gitlab.
    \end{block}
    \begin{center}
        
\includegraphics[width=3in]{images/forgejo-logo.png}
    \end{center}

    Vamos a usar un servidor de Forgejo hosteado en \href{https://git.cubawiki.com.ar/}{git.cubawiki.com.ar}.

    Hay otros proyectos que comparten los \textit{commits} por mail, por ejemplo.
    % si hay una fotito de esto puede quedar bueno
\end{frame}


\begin{frame}[t]{Obteniendo un repositorio Git de otro lado}
    \begin{comando}
        git clone
    \end{comando}

    \begin{block}{}
        Para obtener una \textbf{copia local} de un repositorio existente en algún servidor,
        utilizamos el comando \texttt{git clone [URL]}.

        En el caso de haber creado el repositorio con \texttt{git init}, esto no es necesario. El comando \texttt{clone} es como hacer \textit{init + remote add + pull}.
    \end{block}
    \begin{block}{}
        Importante: \texttt{git clone [URL]} crea un nuevo directorio con el
        repositorio clonado. Debemos hacer \texttt{cd <repo>} luego de clonarlo!
    \end{block}
\end{frame}
